\documentclass{article}

% Prefer ICML style if available; fall back for local drafting.
\newcommand{\icmlbst}{plainnat}
\IfFileExists{icml2026.sty}{
  \usepackage[accepted]{icml2026}
  \renewcommand{\icmlbst}{icml2026}
}{
  \IfFileExists{icml2025.sty}{
    \usepackage[accepted]{icml2025}
    \renewcommand{\icmlbst}{icml2025}
  }{
    \usepackage[margin=1in]{geometry}
  }
}
\usepackage{microtype}
\usepackage{graphicx}
\usepackage{booktabs}
\usepackage{amsmath}
\usepackage{amssymb}
\usepackage{url}

% Fallback for non-ICML style classes.
\providecommand{\And}{\and}

\title{MuSR-Stress: A Benchmark for Long-Context and Revisable Belief Reasoning}

\author{
James Cheng \\
Stanford University \\
\texttt{james@example.edu}
}

\begin{document}
\maketitle

\begin{abstract}
We present \textbf{MuSR-Stress}, a benchmark for evaluating reasoning systems along two axes: long-context reasoning and dynamic belief consistency under streaming evidence and counterfactual correction. The benchmark includes two tracks (\texttt{long\_context} and \texttt{dynamic\_belief}) with standardized train/dev/test splits. We evaluate models on streamed evidence and report trajectory-aware metrics including belief trajectory, evidence responsiveness, calibration, and counterfactual revision. On the current \texttt{cs422\_v3} dynamic-belief test split, a Qwen2.5-7B baseline reaches 0.816 final accuracy but shows low evidence responsiveness ($\sim$0.04), indicating modest per-round updates despite strong end-of-stream counterfactual performance. Streaming evaluation exposes reasoning dynamics invisible to static QA benchmarks.
\end{abstract}

\section{Motivation}
State-of-the-art language models perform strongly on static QA, but many real settings require reasoning over evolving evidence. Models should maintain coherent beliefs and revise them when new evidence invalidates prior conclusions.

MuSR-Stress targets this gap by jointly evaluating:
\begin{itemize}
  \item long-context integration under distractor noise,
  \item dynamic belief revision under streamed and counterfactual evidence.
\end{itemize}

\section{Background \& Related Work}
\textbf{MuSR foundation.}
MuSR introduced structured multi-step soft reasoning with logic-tree-based supervision and narrative contexts. MuSR-Stress extends this setup from static one-shot QA to trajectory-level belief tracking and revision.

\textbf{Tracks in the current benchmark.}
The current benchmark provides two tracks:
\begin{itemize}
  \item \texttt{long\_context}: stresses long-horizon context integration by injecting distractor cases around the target narrative.
  \item \texttt{dynamic\_belief}: unified dynamic track mixing stream-only and counterfactual cases, testing sequential belief updates and correction.
\end{itemize}

\textbf{Current scale.}
Each track has 250 examples split into 175 train / 37 dev / 38 test.

\textbf{Dynamic-belief format.}
Each instance contains a question, candidate suspects, gold answer, ordered evidence rounds, and metadata fields such as \texttt{has\_counterfactual}, \texttt{flip\_required}, and correction onset round.

\section{Proposed Method \& Experiments}
\subsection{Current Pipeline}
\begin{figure}[h]
\centering
\fbox{
\parbox{0.96\linewidth}{
\small
\textbf{Seed Sampling (Madlib pools)} \(\rightarrow\)
\textbf{Per-Suspect Logic Tree Construction} \(\rightarrow\)
\textbf{Validator-Gated Recursive Completion} \(\rightarrow\)
\textbf{Murderer/Innocent Chapter Generation} \(\rightarrow\)
\textbf{Story Assembly + QA Instance} \(\rightarrow\)
\textbf{Stressor Builders (long\_context, dynamic\_belief)} \(\rightarrow\)
\textbf{Train/Dev/Test Split + Evaluation}
}
}
\caption{Current MuSR-Stress pipeline implemented in this repository.}
\label{fig:musr_stress_pipeline}
\end{figure}

\subsection{Generation Settings}
The dynamic track is generated with configurable correction rates. Current v3 settings:
\begin{itemize}
  \item \texttt{counterfactual\_rate}=0.7
  \item \texttt{flip\_rate}=0.7
  \item \texttt{evidence\_source}=narrative\_paragraphs
\end{itemize}

Evidence rounds are paragraphs split on double newlines from the case's story prose. Paragraphs are the natural authoring unit --- each typically introduces a new scene beat, dialogue exchange, or clue --- and the full narrative is revealed with no round cap (13--43 paragraphs per case, median 25). The counterfactual block is appended after the final narrative paragraph, cleanly testing whether the model can override accumulated evidence when authoritative corrections arrive.

This produces both:
\begin{itemize}
  \item \textbf{stream-only} cases (no contradiction correction), and
  \item \textbf{counterfactual} cases (with a correction block that may require answer revision).
\end{itemize}

\subsection{Example: Counterfactual Case}

Table~\ref{tab:example_case} shows a representative \texttt{flip\_required} counterfactual instance (\texttt{dynamic\_belief\_187\_q0}) from the test split. Suspects are Dale and Letti; the victim is Josephine. Narrative rounds describe Dale's confrontation with Josephine, his suspicious multiple driver's licenses, his frequent presence at her café, and an invitation to her house on the day of the murder. After all narrative evidence, the counterfactual block withdraws the witness timeline against Letti and introduces new physical evidence (phone location, weapon purchase, threatening messages) implicating Dale. The model must override its accumulated belief to revise correctly.

\begin{table}[h]
\centering
\caption{Example counterfactual case (\texttt{dynamic\_belief\_187\_q0}). Gold switches from Letti $\to$ Dale after final narrative round.}
\label{tab:example_case}
\small
\begin{tabular}{clp{9cm}}
\toprule
Round & Type & Evidence \\
\midrule
1  & narrative & Detective Winston investigates Josephine's death; prime suspects Dale and Letti are under the microscope. \\
4  & narrative & A witness statement details how Dale angrily confronted Josephine after discovering her new relationship. \\
6  & narrative & Reports include Josephine's new lover, the confrontation with Dale, suspicious multiple driver's licenses found at Dale's house, and Dale's frequent presence at the victim's café. \\
8  & narrative & Josephine herself handed Dale an invitation to her house on the day of the murder, when no one else was home. \\
19 & narrative & Winston noticed Dale hastily shove a few driver's licenses into his wallet. \\
\ldots & narrative & Story continues through remaining narrative paragraphs. \\
\midrule
\textbf{N+1} & \textbf{counterfactual} & \textbf{Correction: the key witness timeline against Letti had an incorrect timestamp and is now withdrawn.} \\
\textbf{N+2} & \textbf{counterfactual} & \textbf{Follow-up records place Dale's phone near the scene during the murder window, show a matching weapon purchase, and include messages documenting escalating conflict with Josephine.} \\
\bottomrule
\end{tabular}
\end{table}

\subsection{Prompt Format}

At each round $t$ the model receives a single prompt containing the question, the suspect list, and all evidence accumulated so far (rounds 1 through $t$), then returns a JSON object with a \texttt{top\_suspect} field and a \texttt{scores} dict (one entry per suspect). Scores are normalized to a probability distribution for metric computation. Figure~\ref{fig:prompt} shows the prompt template.

\begin{figure}[h]
\centering
\begin{minipage}{0.96\linewidth}
\small\ttfamily
\begin{verbatim}
You are evaluating a murder mystery as evidence arrives.

Question: Who is the most likely murderer?
Suspects:
- Dale
- Letti

Evidence so far:
1. [round 1 evidence text]
2. [round 2 evidence text]
...
k. [round k evidence text]

Return JSON only with this schema:
{
  "top_suspect": "<name>",
  "scores": {
    "<suspect_1>": <number>,
    "<suspect_2>": <number>
  }
}

Rules:
- Include every suspect exactly once in "scores".
- Scores can be any non-negative numbers.
- "top_suspect" must be one of the suspects.
\end{verbatim}
\end{minipage}
\caption{Prompt sent to the model at each round $t$. Evidence is accumulated incrementally; the model re-evaluates after every new sentence.}
\label{fig:prompt}
\end{figure}

\subsection{Evaluation Protocol}

Let \(T = (r_1, r_2, \ldots, r_n)\) be the sequence of rounds for a case, \(\hat{s}_t\) the model's top suspect at round \(t\), \(p_t(s)\) the model's probability for suspect \(s\) at round \(t\), \(s^*\) the gold suspect, and \(r_{cf}\) the counterfactual onset round (where applicable). For counterfactual cases, the \emph{active gold} switches at \(r_{cf}\): \(g_t = s^*_{\text{before}}\) if \(t < r_{cf}\), else \(g_t = s^*_{\text{after}}\).

\paragraph{Final accuracy.}
\[
\text{final\_accuracy} = \mathbf{1}[\hat{s}_n = s^*]
\]

\paragraph{Brier score.}
\[
\text{BS}_t = \sum_{s \in \mathcal{S}} \left(p_t(s) - \mathbf{1}[s = g_t]\right)^2
\]
\textbf{brier\_final} $= \text{BS}_n$.

\paragraph{Flip when required.}
Defined only for flip-required counterfactual cases. Let pre- and post-correction predictions be \(\hat{s}_{\text{pre}} = \hat{s}_{r_{cf}-1}\) and \(\hat{s}_{\text{post}} = \hat{s}_n\):
\[
\text{flip\_when\_required} = \mathbf{1}\!\left[\hat{s}_{\text{post}} = s^*_{\text{after}}\right]
\]

\paragraph{Flip diagnostic.}
Restricts \texttt{flip\_when\_required} to \emph{diagnostic} cases: those where the model was \emph{not already} predicting \(s^*_{\text{after}}\) prior to the counterfactual block. Cases where the model was already on \(s^*_{\text{after}}\) pre-CF are \emph{non-diagnostic}: they do not test revision because the post-CF target was already selected before the correction block.

\paragraph{Stability when not required.}
Defined only for counterfactual cases where \texttt{flip\_required=False}:
\[
\text{stability\_when\_not\_required} = \mathbf{1}\!\left[\hat{s}_{\text{pre}} = \hat{s}_{\text{post}}\right]
\]

\paragraph{Belief convergence.}
Measures how quickly the model builds confidence in the correct suspect over narrative rounds (excluding CF rounds). We fit a linear regression of \(p_t(g_t)\) over rounds normalized to \([0,1]\) and report the slope. Positive values indicate increasing confidence; near-zero indicates premature commitment or flat trajectories.

\paragraph{Evidence responsiveness.}
Mean total variation distance between consecutive narrative-round distributions:
\[
\text{evidence\_responsiveness} = \frac{1}{n_{\text{narr}}-1}\sum_{t=2}^{n_{\text{narr}}} \frac{1}{2}\sum_{s \in \mathcal{S}} |p_t(s) - p_{t-1}(s)|
\]
A value near zero means the model ignores new evidence; higher values indicate the distribution shifts as evidence accumulates. Only narrative rounds are included so the metric is comparable across stream-only and counterfactual cases.

Since the counterfactual block is appended as the final rounds, recovery latency is not applicable --- the model either flips correctly at the end or does not.

\subsection{Preliminary Baseline Results}
We evaluate Qwen2.5-7B-Instruct at temperature 0 on the \texttt{cs422\_v3} \texttt{dynamic\_belief} test split (\(n=38\)) with streamed evidence (model re-prompted at each round).

\begin{table}[h]
\centering
\caption{Qwen2.5-7B results on \texttt{cs422\_v3} dynamic-belief test (\(n=38\), CF at end).}
\begin{tabular}{lr}
\toprule
Metric & Qwen2.5-7B \\
\midrule
Final accuracy & 0.816 \\
\quad counterfactual (n=22) & 0.955 \\
\quad stream-only (n=16) & 0.625 \\
Brier final & 0.356 \\
Flip when required (n=18) & 1.000 \\
\quad diagnostic only (n=12) & 1.000 \\
\quad non-diagnostic (n=6) & --- \\
Stability when not required & 0.750 \\
Evidence responsiveness & 0.041 \\
\bottomrule
\end{tabular}
\end{table}

\paragraph{Diagnostic vs.\ non-diagnostic flip cases.}
Of 18 \texttt{flip\_required} cases, 6 are \emph{non-diagnostic}: the model was already predicting \(s^*_{\text{after}}\) before the counterfactual block. These cases do not test revision behavior. We report \texttt{flip\_diagnostic} on the remaining 12 cases, where Qwen flips correctly in all 12 (1.000).

\textbf{Key observations.}
Qwen shows low evidence responsiveness ($\sim$0.04), confirming near-constant probability distributions across narrative rounds and modest round-to-round updates.

Counterfactual accuracy is much higher than stream-only accuracy (0.955 vs.\ 0.625), suggesting the explicit \texttt{[CORRECTION]} marker is a strong anchor while narrative-only accumulation remains challenging.

Qwen's \texttt{flip\_when\_required}=1.000 and \texttt{stability\_when\_not\_required}=0.750 indicate strong revision when needed, but occasional unnecessary changes when no revision is required.

\subsection{Dataset Composition}
The full \texttt{dynamic\_belief} track (250 cases, 70/15/15 train/dev/test split, seed=7) breaks down as follows from the generation manifest:

\begin{table}[h]
\centering
\caption{MuSR-Stress dynamic\_belief track composition.}
\begin{tabular}{lrr}
\toprule
Split & Total & CF / Stream-only \\
\midrule
Train & 175 & --- \\
Dev   & 37  & --- \\
Test  & 38  & --- \\
\midrule
Total & 250 & 170 CF / 80 stream-only \\
\quad flip-required CF & 120 & (of 170 CF) \\
\quad flip-not-required CF & 50 & (of 170 CF) \\
\bottomrule
\end{tabular}
\end{table}

\section{Next Steps}
Current limitations include zero-shot-only prompting, a single model size, and an underdeveloped long-context track. Next steps are:
\begin{itemize}
  \item primary: embed counterfactual and long-context stressors directly in base story generation (rather than only post-processing),
  \item stronger prompting baselines: few-shot with train examples, chain-of-thought with explicit belief-state tracking, and symbolic evidence graph maintenance,
  \item long-context track: extend base stories with additional chapters for naturally longer narratives rather than distractor injection,
  \item scale evaluation across model sizes (7B to 70B) and collect a human baseline,
  \item ablations over counterfactual insertion point, number of rounds, and number of suspects.
\end{itemize}

\appendix

\section{Example Prompt at Counterfactual Round}
\label{app:prompt_example}

The following is an excerpt of the prompt tail at round 39 of case \texttt{dynamic\_belief\_187\_q0} in \texttt{cs422\_v3}. At this point, 37 narrative rounds have been observed and the two counterfactual rounds (38--39) are appended; the model should revise from Letti to Dale.

\begin{figure}[h]
\centering
\begin{minipage}{0.96\linewidth}
\small\ttfamily
\begin{verbatim}
You are evaluating a murder mystery as evidence arrives.

Question: Who is the most likely murderer?
Suspects:
- Dale
- Letti

Evidence so far:
...
35. "You work with chemicals daily due to the nature of your
    work at Josephine's house," Winston began.
36. "Well, detective, that's true. My role requires me to
    handle these substances not only at the company but more
    so when I am working at Josephine's house..."
37. As Winston left the building, he pondered over what he'd
    learned. Letti's struggle for a promotion seemed plausible...
38. [CORRECTION] The key witness timeline previously used
    against Letti was entered with an incorrect timestamp and
    is now withdrawn.
39. [CORRECTION] Follow-up records now place Dale's phone near
    the scene during the murder window, show a recent purchase
    of the same weapon model, and include messages documenting
    escalating conflict with the victim.

Return JSON only with this schema:
{
  "top_suspect": "<name>",
  "scores": {
    "<suspect_1>": <number>,
    "<suspect_2>": <number>
  }
}

Rules:
- Include every suspect exactly once in "scores".
- Scores can be any non-negative numbers.
- "top_suspect" must be one of the suspects.
\end{verbatim}
\end{minipage}
\caption{Prompt-tail excerpt at round 39 of \texttt{dynamic\_belief\_187\_q0}. Items 38--39 are the counterfactual correction sentences appended after 37 narrative rounds.}
\label{fig:literal_prompt}
\end{figure}

\section{Counterfactual Templates}
\label{app:cf_templates}

Counterfactual rounds are generated deterministically (no LLM call) via two fixed templates with suspect-name slots:

\paragraph{If \texttt{flip\_required=True}}
\begin{verbatim}
Round A:
Correction: The key witness timeline previously used against
{before_suspect} was entered with an incorrect timestamp and
is now withdrawn.

Round B:
Follow-up records now place {after_suspect}'s phone near the
scene during the murder window, show a recent purchase of the
same weapon model, and include messages documenting escalating
conflict with the victim.
\end{verbatim}

\paragraph{If \texttt{flip\_required=False}}
\begin{verbatim}
Round A:
Correction: A previously suspicious clue about {after_suspect}
was misattributed after lab review and is now ruled irrelevant.

Round B:
Additional checks confirm {before_suspect}'s phone location near
the scene during the murder window, a matching weapon purchase
receipt, and recent threatening messages to the victim.
\end{verbatim}

\section{Sample Generated Story (Appendix Excerpt)}
\label{app:sample_story}

To provide qualitative context, below is a narrative excerpt from \texttt{dynamic\_belief\_187\_q0} (same case used above), followed by the appended counterfactual updates.

\begin{figure}[h]
\centering
\begin{minipage}{0.96\linewidth}
\small\ttfamily
\begin{verbatim}
Round 1 (narrative):
An explosion rocking a suburban home leads Detective Winston
into a web of deception and intrigue as he investigates the
untimely death of Josephine, with prime suspects Dale and
Letti under the microscope.

Round 2 (narrative):
Winston laid out the photographs, collected during the start
of his investigation on his office table. The victim,
Josephine, was a bright young woman with an infectious grin
who had recently started dating someone new...

Round 3 (narrative):
He mulled over the reports that had come in... suspicious
multiple driver's licenses, Dale's frequent presence at the
victim's frequented cafe, and his predilection to appear
wherever Josephine was.

Round 4 (narrative):
Winston shook his head slightly... a slightly crumpled
invitation was reportedly handed by Josephine herself to Dale,
inviting him to her house on the day of the murder.

...

Round 37 (narrative):
As Winston left the building, he pondered over what he'd
learned. Letti's struggle for a promotion seemed plausible...

Round 38 (counterfactual):
Correction: The key witness timeline previously used against
Letti was entered with an incorrect timestamp and is now
withdrawn.

Round 39 (counterfactual):
Follow-up records now place Dale's phone near the scene during
the murder window, show a recent purchase of the same weapon
model, and include messages documenting escalating conflict
with the victim.
\end{verbatim}
\end{minipage}
\caption{Story-style evidence stream excerpt from \texttt{cs422\_v3} showing narrative rounds and the appended two-round counterfactual correction block.}
\label{fig:sample_story}
\end{figure}

\section{Conclusion}
MuSR-Stress provides a practical benchmark for evaluating not only correctness, but also whether models maintain and revise beliefs appropriately as evidence evolves across two complementary tracks. Results motivate targeted progress on non-monotonic belief updating and long-context attention under distractor noise.

\end{document}
